% ============================================================
% fig_tech_route.tex  七阶段技术路线流程图（Beamer版）
% 依赖：fig_colors.tex 中定义的颜色
% 用法：% ============================================================
% fig_tech_route.tex  七阶段技术路线流程图（Beamer版）
% 依赖：fig_colors.tex 中定义的颜色
% 用法：% ============================================================
% fig_tech_route.tex  七阶段技术路线流程图（Beamer版）
% 依赖：fig_colors.tex 中定义的颜色
% 用法：% ============================================================
% fig_tech_route.tex  七阶段技术路线流程图（Beamer版）
% 依赖：fig_colors.tex 中定义的颜色
% 用法：\input{figures/fig_tech_route.tex}
% 设计尺寸：约 9.5cm × 5.6cm，适合单页 Beamer 帧（居中使用）
% ============================================================
\begin{tikzpicture}[
  node distance=0.15cm,
  block/.style={rectangle, draw=deepblue, fill=skyblue,
                text width=9.5cm, align=center,
                rounded corners=3pt, minimum height=0.55cm,
                font=\footnotesize},
  coreblock/.style={rectangle, draw=deeppink, fill=improvedpink,
                    text width=9.5cm, align=center,
                    rounded corners=3pt, minimum height=0.55cm,
                    font=\footnotesize\bfseries},
  endblock/.style={rectangle, draw=deepblue, fill=outputpink,
                   text width=9.5cm, align=center,
                   rounded corners=3pt, minimum height=0.55cm,
                   font=\footnotesize},
  arrow/.style={thick, -Stealth, deepblue}
]

\node (s1) [block]
  {\textbf{阶段一：文献调研与资料收集}\enspace
   阅读 Apollo 官方文档与源码；收集路径规划相关学术论文};

\node (s2) [block, below=of s1]
  {\textbf{阶段二：Apollo 源码分析与理解}\enspace
   Planning 模块解析；核心数据结构与算法流程梳理};

\node (s3) [block, below=of s2]
  {\textbf{阶段三：问题识别与定量分析}\enspace
   多障碍物场景性能瓶颈定位；ST 边界与 PathDecider 问题分析};

\node (s4) [coreblock, below=of s3]
  {\textbf{阶段四：改进算法设计（核心）}\enspace
   威胁评估模型；时空可通行走廊算法；Apollo 集成方案};

\node (s5) [coreblock, below=of s4]
  {\textbf{阶段五：算法实现与调试（核心）}\enspace
   代码开发与模块集成；参数调优；单元与功能测试};

\node (s6) [block, below=of s5]
  {\textbf{阶段六：实验验证与分析}\enspace
   仿真场景设计与数据采集；性能指标对比与结果分析};

\node (s7) [endblock, below=of s6]
  {\textbf{阶段七：论文撰写与答辩}};

\draw[arrow] (s1) -- (s2);
\draw[arrow] (s2) -- (s3);
\draw[arrow] (s3) -- (s4);
\draw[arrow] (s4) -- (s5);
\draw[arrow] (s5) -- (s6);
\draw[arrow] (s6) -- (s7);

\end{tikzpicture}

% 设计尺寸：约 9.5cm × 5.6cm，适合单页 Beamer 帧（居中使用）
% ============================================================
\begin{tikzpicture}[
  node distance=0.15cm,
  block/.style={rectangle, draw=deepblue, fill=skyblue,
                text width=9.5cm, align=center,
                rounded corners=3pt, minimum height=0.55cm,
                font=\footnotesize},
  coreblock/.style={rectangle, draw=deeppink, fill=improvedpink,
                    text width=9.5cm, align=center,
                    rounded corners=3pt, minimum height=0.55cm,
                    font=\footnotesize\bfseries},
  endblock/.style={rectangle, draw=deepblue, fill=outputpink,
                   text width=9.5cm, align=center,
                   rounded corners=3pt, minimum height=0.55cm,
                   font=\footnotesize},
  arrow/.style={thick, -Stealth, deepblue}
]

\node (s1) [block]
  {\textbf{阶段一：文献调研与资料收集}\enspace
   阅读 Apollo 官方文档与源码；收集路径规划相关学术论文};

\node (s2) [block, below=of s1]
  {\textbf{阶段二：Apollo 源码分析与理解}\enspace
   Planning 模块解析；核心数据结构与算法流程梳理};

\node (s3) [block, below=of s2]
  {\textbf{阶段三：问题识别与定量分析}\enspace
   多障碍物场景性能瓶颈定位；ST 边界与 PathDecider 问题分析};

\node (s4) [coreblock, below=of s3]
  {\textbf{阶段四：改进算法设计（核心）}\enspace
   威胁评估模型；时空可通行走廊算法；Apollo 集成方案};

\node (s5) [coreblock, below=of s4]
  {\textbf{阶段五：算法实现与调试（核心）}\enspace
   代码开发与模块集成；参数调优；单元与功能测试};

\node (s6) [block, below=of s5]
  {\textbf{阶段六：实验验证与分析}\enspace
   仿真场景设计与数据采集；性能指标对比与结果分析};

\node (s7) [endblock, below=of s6]
  {\textbf{阶段七：论文撰写与答辩}};

\draw[arrow] (s1) -- (s2);
\draw[arrow] (s2) -- (s3);
\draw[arrow] (s3) -- (s4);
\draw[arrow] (s4) -- (s5);
\draw[arrow] (s5) -- (s6);
\draw[arrow] (s6) -- (s7);

\end{tikzpicture}

% 设计尺寸：约 9.5cm × 5.6cm，适合单页 Beamer 帧（居中使用）
% ============================================================
\begin{tikzpicture}[
  node distance=0.15cm,
  block/.style={rectangle, draw=deepblue, fill=skyblue,
                text width=9.5cm, align=center,
                rounded corners=3pt, minimum height=0.55cm,
                font=\footnotesize},
  coreblock/.style={rectangle, draw=deeppink, fill=improvedpink,
                    text width=9.5cm, align=center,
                    rounded corners=3pt, minimum height=0.55cm,
                    font=\footnotesize\bfseries},
  endblock/.style={rectangle, draw=deepblue, fill=outputpink,
                   text width=9.5cm, align=center,
                   rounded corners=3pt, minimum height=0.55cm,
                   font=\footnotesize},
  arrow/.style={thick, -Stealth, deepblue}
]

\node (s1) [block]
  {\textbf{阶段一：文献调研与资料收集}\enspace
   阅读 Apollo 官方文档与源码；收集路径规划相关学术论文};

\node (s2) [block, below=of s1]
  {\textbf{阶段二：Apollo 源码分析与理解}\enspace
   Planning 模块解析；核心数据结构与算法流程梳理};

\node (s3) [block, below=of s2]
  {\textbf{阶段三：问题识别与定量分析}\enspace
   多障碍物场景性能瓶颈定位；ST 边界与 PathDecider 问题分析};

\node (s4) [coreblock, below=of s3]
  {\textbf{阶段四：改进算法设计（核心）}\enspace
   威胁评估模型；时空可通行走廊算法；Apollo 集成方案};

\node (s5) [coreblock, below=of s4]
  {\textbf{阶段五：算法实现与调试（核心）}\enspace
   代码开发与模块集成；参数调优；单元与功能测试};

\node (s6) [block, below=of s5]
  {\textbf{阶段六：实验验证与分析}\enspace
   仿真场景设计与数据采集；性能指标对比与结果分析};

\node (s7) [endblock, below=of s6]
  {\textbf{阶段七：论文撰写与答辩}};

\draw[arrow] (s1) -- (s2);
\draw[arrow] (s2) -- (s3);
\draw[arrow] (s3) -- (s4);
\draw[arrow] (s4) -- (s5);
\draw[arrow] (s5) -- (s6);
\draw[arrow] (s6) -- (s7);

\end{tikzpicture}

% 设计尺寸：约 9.5cm × 5.6cm，适合单页 Beamer 帧（居中使用）
% ============================================================
\begin{tikzpicture}[
  node distance=0.15cm,
  block/.style={rectangle, draw=deepblue, fill=skyblue,
                text width=9.5cm, align=center,
                rounded corners=3pt, minimum height=0.55cm,
                font=\footnotesize},
  coreblock/.style={rectangle, draw=deeppink, fill=improvedpink,
                    text width=9.5cm, align=center,
                    rounded corners=3pt, minimum height=0.55cm,
                    font=\footnotesize\bfseries},
  endblock/.style={rectangle, draw=deepblue, fill=outputpink,
                   text width=9.5cm, align=center,
                   rounded corners=3pt, minimum height=0.55cm,
                   font=\footnotesize},
  arrow/.style={thick, -Stealth, deepblue}
]

\node (s1) [block]
  {\textbf{阶段一：文献调研与资料收集}\enspace
   阅读 Apollo 官方文档与源码；收集路径规划相关学术论文};

\node (s2) [block, below=of s1]
  {\textbf{阶段二：Apollo 源码分析与理解}\enspace
   Planning 模块解析；核心数据结构与算法流程梳理};

\node (s3) [block, below=of s2]
  {\textbf{阶段三：问题识别与定量分析}\enspace
   多障碍物场景性能瓶颈定位；ST 边界与 PathDecider 问题分析};

\node (s4) [coreblock, below=of s3]
  {\textbf{阶段四：改进算法设计（核心）}\enspace
   威胁评估模型；时空可通行走廊算法；Apollo 集成方案};

\node (s5) [coreblock, below=of s4]
  {\textbf{阶段五：算法实现与调试（核心）}\enspace
   代码开发与模块集成；参数调优；单元与功能测试};

\node (s6) [block, below=of s5]
  {\textbf{阶段六：实验验证与分析}\enspace
   仿真场景设计与数据采集；性能指标对比与结果分析};

\node (s7) [endblock, below=of s6]
  {\textbf{阶段七：论文撰写与答辩}};

\draw[arrow] (s1) -- (s2);
\draw[arrow] (s2) -- (s3);
\draw[arrow] (s3) -- (s4);
\draw[arrow] (s4) -- (s5);
\draw[arrow] (s5) -- (s6);
\draw[arrow] (s6) -- (s7);

\end{tikzpicture}
